% Supplementary Table S1: included English-language stroke guidelines only.
% Assessment scales and computational mappings are reported separately in Tables S2 and S3.
% Required packages: booktabs, longtable, array, ragged2e, pdflscape.

\begin{landscape}
\clearpage
\begingroup
\renewcommand{\tablename}{Supplementary Table}
\renewcommand{\thetable}{S1}
\setlength{\tabcolsep}{3pt}
\setlength{\LTpre}{0pt}
\setlength{\LTpost}{0pt}
\renewcommand{\arraystretch}{1.08}
\footnotesize
\begin{longtable}{@{}
>{\RaggedRight\arraybackslash}p{2.10cm}
>{\RaggedRight\arraybackslash}p{8.70cm}
>{\RaggedRight\arraybackslash}p{4.20cm}
>{\RaggedRight\arraybackslash}p{8.70cm}
@{}}
\caption{Included English-language stroke guideline documents and their clinical scope.}
\label{tab:s1_guidelines_compact}\\
\toprule
\textbf{Region} & \textbf{Guideline document and year} & \textbf{Issuing organization} & \textbf{Clinical scope} \\
\midrule
\endfirsthead
\multicolumn{4}{l}{\textit{Supplementary Table S1 continued}}\\
\toprule
\textbf{Region} & \textbf{Guideline document and year} & \textbf{Issuing organization} & \textbf{Clinical scope} \\
\midrule
\endhead
\midrule
\multicolumn{4}{r}{\textit{Continued on the next page}}\\
\endfoot
\bottomrule
\endlastfoot
United States & Guidelines for Adult Stroke Rehabilitation and Recovery (2016) \cite{S1Winstein2016AHAASARehab} & AHA/ASA & Cross-setting rehabilitation: assessment, prognosis, intervention and community reintegration \\
United States & VA/DoD Clinical Practice Guideline for Stroke Rehabilitation (2024) \cite{S1VADoD2024StrokeRehab} & VA/DoD & Post-stroke rehabilitation: assessment, prognosis, intervention, transitions and participation \\
United States & Clinical Practice Guideline for AFO and FES Post-Stroke (2021) \cite{S1Johnston2021AFOFESPostStroke} & ANPT/APTA & AFO/FES selection for walking, balance and mobility \\
United States & Clinical Practice Guideline to Improve Locomotor Function (2020) \cite{S1Hornby2020LocomotorCPG} & APTA/ANPT & Intensive task-specific locomotor training in chronic stroke \\
United States & Guideline for Early Management of Acute Ischemic Stroke (2026) \cite{S1Prabhakaran2026AHAASAAIS} & AHA/ASA & Acute ischaemic stroke diagnosis, reperfusion, management and early prevention \\
United States & Guideline for Aneurysmal Subarachnoid Hemorrhage (2023) \cite{S1Hoh2023AHAASAASAH} & AHA/ASA & aSAH diagnosis, acute treatment, complication management and recovery \\
United States & Guideline for Secondary Stroke Prevention (2021) \cite{S1Kleindorfer2021AHAASASecondaryPrevention} & AHA/ASA & Secondary prevention: risk assessment, medication, lifestyle and follow-up \\
United States & Guideline for Spontaneous Intracerebral Hemorrhage (2022) \cite{S1Greenberg2022AHAASAICH} & AHA/ASA & ICH diagnosis, haemostasis, blood-pressure management, surgery and recovery \\
United Kingdom & National Clinical Guideline for Stroke for the UK and Ireland (2023) \cite{S1Intercollegiate2023NationalStroke} & Intercollegiate Stroke Working Party & Stroke services, acute care, rehabilitation and life after stroke \\
United Kingdom & Stroke Rehabilitation in Adults, NICE NG236 (2023) \cite{S1NICE2023NG236} & NICE & Adult rehabilitation: assessment, goals, therapy and community care \\
United Kingdom & Stroke and TIA in Over 16s, NICE NG128 (2019; updated 2022) \cite{S1NICE2019NG128} & NICE & Acute stroke/TIA recognition, imaging, reperfusion and initial management \\
Canada & CSBPR Rehabilitation, Part One: Planning (2025) \cite{S1CSBPR2025RehabPart1} & Heart and Stroke Foundation / CSBPR & Rehabilitation planning: screening, assessment, prognosis, goals and transitions \\
Canada & CSBPR Rehabilitation, Part Two: Delivery (2025) \cite{S1CSBPR2025RehabPart2} & Heart and Stroke Foundation / CSBPR & Rehabilitation delivery: therapy dose, mobility, upper-limb function and cognition \\
Canada & CSBPR Rehabilitation, Part Three: Activity and Community Participation (2025) \cite{S1CSBPR2025RehabPart3} & Heart and Stroke Foundation / CSBPR & ADL, self-management, participation, social roles and long-term follow-up \\
Canada & CSBPR Acute Stroke Management (2022; EVT update 2025) \cite{S1CSBPR20222025Acute} & Heart and Stroke Foundation / CSBPR & Acute management: recognition, imaging, reperfusion and stroke-unit care \\
Canada & CSBPR Management of Spontaneous ICH (2020) \cite{S1CSBPR2020ICH} & Heart and Stroke Foundation / CSBPR & Acute ICH management, complications and recovery planning \\
Canada & CSBPR Secondary Prevention of Stroke (2020) \cite{S1CSBPR2020SecondaryPrevention} & Heart and Stroke Foundation / CSBPR & Secondary prevention: vascular risk, antithrombotics and lifestyle \\
Australia/New Zealand & Australian and New Zealand Living Clinical Guidelines, Chapter 5: Rehabilitation (2025) \cite{S1StrokeFoundationLivingGuidelineRehab} & Stroke Foundation & Rehabilitation assessment and intervention across settings, including participation \\
Australia/New Zealand & National Rehabilitation Stroke Services Framework (2022) \cite{S1StrokeFoundation2022RehabFramework} & Stroke Foundation & Rehabilitation services, workforce, reintegration and quality improvement \\
Australia/New Zealand & Acute Stroke Clinical Care Standard (2019) \cite{S1ACSQHC2019AcuteStrokeStandard} & ACSQHC & Acute-care standard: recognition, treatment, inpatient care and rehabilitation transition \\
Europe & ESO Guideline on Stroke Rehabilitation: Motor Rehabilitation (2025) \cite{S1AltMurphy2025ESOMotorRehab} & ESO & Post-stroke motor-rehabilitation assessment and intervention \\
Europe & EAN/EFNS Guideline on Pharmacological Support in Early Motor Rehabilitation (2021) \cite{S1Beghi2021EANEFNSPharmacological} & EAN/EFNS & Pharmacological support for early motor rehabilitation after ischaemic stroke \\
South Korea & Clinical Practice Guideline for Stroke Rehabilitation in Korea, Part 1 (2022; published 2023) \cite{S1Korea2022Part1} & Korean Society for NeuroRehabilitation & Upper- and lower-limb function, gait, balance and task practice \\
South Korea & Clinical Practice Guideline for Stroke Rehabilitation in Korea, Part 2 (2024; published 2025) \cite{S1Lee2025KoreaPart2} & Korean Society for NeuroRehabilitation & Motor and sensory rehabilitation, upper limb, assistive technology and shoulder care \\
Mainland China & Chinese Stroke Association Guidelines for Clinical Management of Ischaemic Cerebrovascular Diseases: Executive Summary and 2023 Update \cite{S1Liu2023CSAIschaemicUpdate} & Chinese Stroke Association & Ischaemic cerebrovascular disease diagnosis, acute treatment, complication management and long-term risk-factor control \\
Singapore & National One-Rehabilitation Framework (2019) \cite{S1MOH2019OneRehabFramework} & Ministry of Health & Integrated rehabilitation pathways, standardized assessment and community reintegration \\
Hong Kong & Strategic Service Framework for Rehabilitation Services (2013) \cite{S1HospitalAuthority2013RehabSSF} & Hospital Authority & Rehabilitation service planning, capacity, continuity and participation \\
Hong Kong & Consensus Statement on Ischaemic Stroke Care in Hong Kong (2004) \cite{S1HongKong2004IschaemicStrokeConsensus} & Hong Kong Neurological/Stroke Societies & Acute ischaemic stroke assessment, treatment and recurrence prevention \\
\end{longtable}
\noindent\footnotesize\textit{Notes.} The table contains 28 documents from nine regions. Eligibility required a complete English full text; an English title or abstract alone was insufficient. Guideline titles and issuing organizations are reported as published, and clinical scope is a concise paraphrase of each document's coverage. Assessment scales identified through the guideline stream are catalogued separately in Supplementary Table~S2. ACSQHC, Australian Commission on Safety and Quality in Health Care; ADL, activities of daily living; AFO, ankle--foot orthosis; AHA/ASA, American Heart Association/American Stroke Association; ANPT/APTA, Academy of Neurologic Physical Therapy/American Physical Therapy Association; aSAH, aneurysmal subarachnoid haemorrhage; CSA, Chinese Stroke Association; CSBPR, Canadian Stroke Best Practice Recommendations; EAN/EFNS, European Academy of Neurology/European Federation of Neurorehabilitation Societies; ESO, European Stroke Organisation; EVT, endovascular therapy; FES, functional electrical stimulation; ICH, intracerebral haemorrhage; NICE, National Institute for Health and Care Excellence; TIA, transient ischaemic attack; VA/DoD, U.S. Department of Veterans Affairs/U.S. Department of Defense.
\endgroup
\end{landscape}
